% !TeX root = semana04.tex
% !TeX program = pdflatex
% Material de la semana 04: sesión de 3 horas después de la evaluación.
\documentclass[10pt,a5paper]{article}
\newif\ifmostrarSoluciones
\mostrarSolucionestrue
\usepackage{tikz}
\usepackage{estilo_apuntes}
\usetikzlibrary{arrows.meta,positioning,babel}
\encabezado{Semana 4}{Apunte semanal}

\begin{document}
\tituloclase{Límites infinitos y al infinito}{Asíntotas verticales y horizontales}{Semana 4 · 3 horas}

\begin{objetivos}
\begin{itemize}
  \item Interpretar límites infinitos y límites al infinito mediante expresiones y gráficas.
  \item Identificar asíntotas verticales y horizontales a partir de límites.
  \item Calcular límites al infinito de funciones racionales mediante comparación de grados.
  \item Determinar el crecimiento no acotado de polinomios y funciones racionales cuando $x\to\pm\infty$.
  \item Reconocer casos en los que un denominador nulo no produce una asíntota vertical y casos que requieren usar $\sqrt{x^2}=|x|$.
\end{itemize}
\end{objetivos}

\begin{resumen}
Un límite infinito describe valores de la función que crecen o decrecen sin cota cuando $x$ se aproxima a un número. Un límite al infinito describe el comportamiento de la función cuando $x$ toma valores positivos o negativos cada vez mayores en valor absoluto. Estas dos ideas permiten identificar asíntotas verticales y horizontales.
\end{resumen}

\begin{nota}[organización de la sesión]
La primera evaluación ocupa dos horas de la semana. Este material está diseñado para las tres horas restantes: aproximadamente 40 minutos para límites infinitos y asíntotas verticales, 35 minutos para límites al infinito y asíntotas horizontales, 40 minutos para comparación de grados, 30 minutos para crecimiento no acotado y el caso especial con raíces, y 35 minutos para práctica y cierre.
\end{nota}

\section*{Dos comportamientos diferentes}

Las expresiones
\[
  \lim_{x\to a}f(x)=+\infty
  \qquad\text{y}\qquad
  \lim_{x\to+\infty}f(x)=L
\]
responden preguntas distintas. En la primera, $x$ se aproxima al número $a$ y los valores de $f(x)$ crecen sin cota. En la segunda, $x$ crece sin cota y los valores de $f(x)$ se aproximan al número real $L$.

\begin{nota}
Los símbolos $+\infty$ y $-\infty$ no representan números reales. Por ello, un límite infinito describe un comportamiento no acotado; no afirma que la función alcance el valor infinito.
\end{nota}

\section*{Límites infinitos y asíntotas verticales}

\begin{definicion}[límite infinito]
Escribimos
\[
  \lim_{x\to a}f(x)=+\infty
\]
cuando los valores de $f(x)$ pueden hacerse positivos y tan grandes como se desee al tomar puntos $x$ del dominio suficientemente próximos a $a$, con $x\neq a$. De manera análoga,
\[
  \lim_{x\to a}f(x)=-\infty
\]
significa que los valores de $f(x)$ pueden hacerse negativos y tan grandes en valor absoluto como se desee.
\end{definicion}

Los límites infinitos también pueden ser laterales. El signo debe analizarse a cada lado del punto, pues una misma función puede tender a $-\infty$ por un lado y a $+\infty$ por el otro.

\begin{definicion}[asíntota vertical]
La recta $x=a$ es una \textbf{asíntota vertical} de la gráfica de $f$ si al menos uno de los límites laterales de $f$ en $a$ es $+\infty$ o $-\infty$.
\end{definicion}

\begin{ejemplo}[signos diferentes a cada lado]
Sea
\[
  f\colon\mathbb{R}\setminus\{2\}\longrightarrow\mathbb{R},
  \qquad
  f(x)=\frac{1}{x-2}.
\]
Determine el comportamiento de $f$ cerca de $2$.
\end{ejemplo}
\ifmostrarSoluciones
\begin{solucion}
Cuando $x<2$, el número $x-2$ es negativo y se aproxima a $0$; por tanto, su recíproco decrece sin cota. Cuando $x>2$, el número $x-2$ es positivo y se aproxima a $0$; por tanto, su recíproco crece sin cota. Así,
\[
  \lim_{x\to2^-}\frac{1}{x-2}=-\infty,
  \qquad
  \lim_{x\to2^+}\frac{1}{x-2}=+\infty.
\]
La recta $x=2$ es una asíntota vertical. Como los comportamientos laterales tienen signos diferentes, no se escribe un único límite bilateral igual a $+\infty$ o a $-\infty$.
\end{solucion}
\else\vspace{20mm}\fi

\begin{center}
\begin{tikzpicture}[x=1.25cm,y=0.72cm,every node/.style={font=\scriptsize}]
  \draw[->,Gris] (-1.2,0)--(5.2,0) node[right] {$x$};
  \draw[->,Gris] (0,-5.2)--(0,5.4) node[above] {$y$};
  \draw[Rojo,dashed,thick] (2,-5)--(2,5) node[above] {$x=2$};
  \draw[Azul,very thick,domain=-1:1.8,samples=120]
    plot (\x,{1/(\x-2)});
  \draw[Azul,very thick,domain=2.2:5,samples=120]
    plot (\x,{1/(\x-2)});
  \node[Azul,right] at (3.3,0.85) {$y=1/(x-2)$};
\end{tikzpicture}
\end{center}

\begin{ejemplo}[el mismo signo a ambos lados]
Calcule los límites laterales de
\[
  g(x)=\frac{1}{(x+1)^2}
\]
cuando $x$ se aproxima a $-1$ e identifique la asíntota vertical.
\end{ejemplo}
\ifmostrarSoluciones
\begin{solucion}
Para $x\neq-1$, el denominador $(x+1)^2$ es positivo. Al aproximarse $x$ a $-1$ por cualquiera de los dos lados, el denominador se aproxima a $0$ mediante valores positivos. En consecuencia,
\[
  \lim_{x\to-1^-}\frac{1}{(x+1)^2}
  =\lim_{x\to-1^+}\frac{1}{(x+1)^2}
  =+\infty.
\]
En este caso también podemos escribir
\[
  \lim_{x\to-1}\frac{1}{(x+1)^2}=+\infty.
\]
La recta $x=-1$ es una asíntota vertical.
\end{solucion}
\else\vspace{18mm}\fi

\begin{nota}[un denominador nulo no es suficiente]
Que el denominador se anule en $x=a$ no garantiza una asíntota vertical. Por ejemplo,
\[
  \frac{x^2-4}{x-2}=x+2
  \qquad (x\neq2).
\]
Por tanto,
\[
  \lim_{x\to2}\frac{x^2-4}{x-2}=4,
\]
y la recta $x=2$ no es una asíntota vertical. Antes de decidir, se deben simplificar los factores comunes y estudiar el límite.
\end{nota}

\section*{Límites al infinito y asíntotas horizontales}

\begin{definicion}[límite al infinito]
Escribimos
\[
  \lim_{x\to+\infty}f(x)=L
\]
cuando los valores de $f(x)$ se aproximan a $L$ al tomar valores positivos de $x$ cada vez mayores. La expresión
\[
  \lim_{x\to-\infty}f(x)=L
\]
describe el comportamiento de $f(x)$ cuando $x$ toma valores negativos cuyo valor absoluto es cada vez mayor.
\end{definicion}

\begin{definicion}[asíntota horizontal]
La recta $y=L$ es una \textbf{asíntota horizontal} de la gráfica de $f$ si se cumple al menos una de las igualdades
\[
  \lim_{x\to+\infty}f(x)=L
  \qquad\text{o}\qquad
  \lim_{x\to-\infty}f(x)=L.
\]
Una función puede tener asíntotas horizontales diferentes hacia $+\infty$ y hacia $-\infty$.
\end{definicion}

\begin{nota}
Una asíntota horizontal describe el comportamiento de la función cuando $x$ se aleja indefinidamente. La gráfica puede cruzar la recta $y=L$ una o varias veces para valores finitos de $x$ y, aun así, aproximarse a ella cuando $x\to+\infty$ o $x\to-\infty$.
\end{nota}

\begin{ejemplo}[una asíntota horizontal]
Sea
\[
  h\colon\mathbb{R}\setminus\{-1\}\longrightarrow\mathbb{R},
  \qquad
  h(x)=\frac{2x+3}{x+1}.
\]
Determine su comportamiento cuando $x\to+\infty$ y cuando $x\to-\infty$.
\end{ejemplo}
\ifmostrarSoluciones
\begin{solucion}
Dividimos el numerador y el denominador entre $x$:
\[
  \frac{2x+3}{x+1}
  =\frac{2+3/x}{1+1/x}.
\]
Como
\[
  \lim_{x\to+\infty}\frac{1}{x}
  =\lim_{x\to-\infty}\frac{1}{x}
  =0,
\]
obtenemos
\[
  \lim_{x\to+\infty}\frac{2x+3}{x+1}
  =\lim_{x\to-\infty}\frac{2x+3}{x+1}
  =2.
\]
La recta $y=2$ es una asíntota horizontal.
\end{solucion}
\else\vspace{18mm}\fi

\begin{center}
\begin{tikzpicture}[x=0.85cm,y=0.85cm,every node/.style={font=\scriptsize}]
  \draw[->,Gris] (-5.2,0)--(5.4,0) node[right] {$x$};
  \draw[->,Gris] (0,-2.0)--(0,5.2) node[above] {$y$};
  \draw[Rojo,dashed,thick] (-5,2)--(5.2,2) node[right] {$y=2$};
  \draw[Rojo,dashed] (-1,-1.8)--(-1,5) node[above] {$x=-1$};
  \draw[Azul,very thick,domain=-5:-1.35,samples=120]
    plot (\x,{(2*\x+3)/(\x+1)});
  \draw[Azul,very thick,domain=-0.65:5,samples=140]
    plot (\x,{(2*\x+3)/(\x+1)});
  \node[Azul,right] at (2.3,2.6) {$y=(2x+3)/(x+1)$};
\end{tikzpicture}
\end{center}

\section*{Comparación de grados}

Considere una función racional
\[
  r(x)=\frac{a_nx^n+\cdots+a_1x+a_0}
             {b_mx^m+\cdots+b_1x+b_0},
  \qquad
  a_n\neq0,
  \quad
  b_m\neq0.
\]
Los términos de mayor grado determinan el comportamiento cuando $x\to+\infty$ o $x\to-\infty$.

\begin{proposicion}[límites al infinito de funciones racionales]
Para la función racional anterior se tienen los siguientes casos:
\begin{enumerate}
  \item Si $n<m$, entonces
  \[
    \lim_{x\to\pm\infty}r(x)=0.
  \]
  La recta $y=0$ es una asíntota horizontal.
  \item Si $n=m$, entonces
  \[
    \lim_{x\to\pm\infty}r(x)=\frac{a_n}{b_m}.
  \]
  La recta $y=a_n/b_m$ es una asíntota horizontal.
  \item Si $n>m$, la función no tiene asíntota horizontal. En este caso,
  \[
    \lim_{x\to\pm\infty}
    \frac{r(x)}{x^{n-m}}
    =\frac{a_n}{b_m},
  \]
  de modo que $r(x)$ tiene el mismo comportamiento dominante que $(a_n/b_m)x^{n-m}$. El signo hacia cada extremo depende del cociente $a_n/b_m$ y de si $n-m$ es par o impar.
\end{enumerate}
\end{proposicion}

\begin{ejemplo}[tres comparaciones de grados]
Calcule los límites al infinito de las siguientes funciones e identifique sus asíntotas horizontales:
\[
  f(x)=\frac{2x+1}{x^2+4},
  \qquad
  g(x)=\frac{3x^2-x+1}{2x^2+5},
  \qquad
  q(x)=\frac{x^3+1}{x^2+1}.
\]
\end{ejemplo}
\ifmostrarSoluciones
\begin{solucion}
En $f$, el grado del numerador es menor que el del denominador; por tanto,
\[
  \lim_{x\to\pm\infty}f(x)=0,
\]
y $y=0$ es una asíntota horizontal. En $g$, los grados son iguales y el cociente de los coeficientes principales es $3/2$; por tanto,
\[
  \lim_{x\to\pm\infty}g(x)=\frac{3}{2},
\]
y $y=3/2$ es una asíntota horizontal. En $q$, el grado del numerador es mayor. Como
\[
  \frac{x^3+1}{x^2+1}
  =x+\frac{1-x}{x^2+1},
\]
la función se comporta como $x$ y no tiene asíntota horizontal.
\end{solucion}
\else\vspace{24mm}\fi

\section*{Crecimiento no acotado cuando $x\to\pm\infty$}

También puede ocurrir que tanto la variable como los valores de la función crezcan sin cota. Por ejemplo,
\[
  \lim_{x\to+\infty}x^2=+\infty.
\]
Esta expresión significa que $x^2$ puede hacerse tan grande como se desee al tomar valores positivos de $x$ suficientemente grandes.

\begin{proposicion}[comportamiento de las potencias]
Sea $k$ un entero positivo. Entonces
\[
  \lim_{x\to+\infty}x^k=+\infty.
\]
Hacia $-\infty$, el resultado depende de la paridad de $k$:
\[
  \lim_{x\to-\infty}x^k=
  \begin{cases}
    +\infty, & k\text{ es par},\\
    -\infty, & k\text{ es impar}.
  \end{cases}
\]
Además,
\[
  \lim_{x\to\pm\infty}\frac{1}{x^k}=0.
\]
\end{proposicion}

Para un polinomio, el término de mayor grado determina el comportamiento en ambos extremos. Para una función racional con mayor grado en el numerador, se compara el cociente con
\[
  \frac{a_n}{b_m}x^{n-m}.
\]

\begin{ejemplo}[signo y paridad del término dominante]
Calcule los límites al infinito de
\[
  f(x)=\frac{3x^3-x}{x^2+1}
  \qquad\text{y}\qquad
  g(x)=\frac{-2x^4+x}{x^2+3}.
\]
\end{ejemplo}
\ifmostrarSoluciones
\begin{solucion}
La función $f$ se comporta como $3x$. Como la potencia dominante es impar y su coeficiente es positivo,
\[
  \lim_{x\to+\infty}f(x)=+\infty,
  \qquad
  \lim_{x\to-\infty}f(x)=-\infty.
\]
La función $g$ se comporta como $-2x^2$. La potencia dominante es par, por lo que $x^2$ es positivo en ambos extremos; el coeficiente negativo determina el signo final:
\[
  \lim_{x\to+\infty}g(x)=-\infty,
  \qquad
  \lim_{x\to-\infty}g(x)=-\infty.
\]
Ninguna de las dos funciones tiene asíntota horizontal.
\end{solucion}
\else\vspace{22mm}\fi

\section*{Caso especial: raíces con $x^2$}

La identidad correcta es
\[
  \sqrt{x^2}=|x|,
\]
no $\sqrt{x^2}=x$ para todo número real $x$. Esta diferencia es esencial cuando se estudian los dos extremos del eje real.

\begin{ejemplo}[dos asíntotas horizontales]
Calcule
\[
  \lim_{x\to+\infty}\frac{\sqrt{x^2+1}}{x}
  \qquad\text{y}\qquad
  \lim_{x\to-\infty}\frac{\sqrt{x^2+1}}{x}.
\]
\end{ejemplo}
\ifmostrarSoluciones
\begin{solucion}
Factorizamos $x^2$ dentro de la raíz:
\[
  \frac{\sqrt{x^2+1}}{x}
  =\frac{\sqrt{x^2\left(1+1/x^2\right)}}{x}
  =\frac{|x|}{x}\sqrt{1+\frac{1}{x^2}}.
\]
Si $x\to+\infty$, entonces $|x|/x=1$; si $x\to-\infty$, entonces $|x|/x=-1$. Por tanto,
\[
  \lim_{x\to+\infty}\frac{\sqrt{x^2+1}}{x}=1,
  \qquad
  \lim_{x\to-\infty}\frac{\sqrt{x^2+1}}{x}=-1.
\]
La gráfica tiene dos asíntotas horizontales: $y=1$ hacia la derecha y $y=-1$ hacia la izquierda.
\end{solucion}
\else\vspace{22mm}\fi

\begin{ejercicios}[práctica de la sesión]
\begin{enumerate}
  \item Calcule los límites laterales de $\displaystyle f(x)=\frac{1}{x+3}$ cuando $x\to-3$ e identifique la asíntota vertical.
  \item Calcule $\displaystyle\lim_{x\to1}\frac{1}{(x-1)^2}$ e interprete el resultado gráficamente.
  \item Decida si $x=2$ es una asíntota vertical de $\displaystyle g(x)=\frac{x^2-4}{x-2}$.
  \item Calcule los límites cuando $x\to\pm\infty$ e identifique las asíntotas horizontales, si existen:
  \begin{enumerate}
    \item $\displaystyle\frac{5x-1}{x^2+2}$,
    \item $\displaystyle\frac{4x^2+3}{2x^2-x}$,
    \item $\displaystyle\frac{x^3-2}{x^2+1}$.
  \end{enumerate}
  \item Determine los límites cuando $x\to+\infty$ y cuando $x\to-\infty$ de
  \[
    p(x)=-3x^5+2x^2-1
    \qquad\text{y}\qquad
    r(x)=\frac{x^4+1}{-x^2+2}.
  \]
  \item Calcule $\displaystyle\lim_{x\to+\infty}\frac{2x+1}{\sqrt{x^2+4}}$ y $\displaystyle\lim_{x\to-\infty}\frac{2x+1}{\sqrt{x^2+4}}$.
  \item Una función satisface
  \[
    \lim_{x\to2^-}f(x)=+\infty,
    \qquad
    \lim_{x\to2^+}f(x)=-\infty,
    \qquad
    \lim_{x\to+\infty}f(x)=3.
  \]
  Indique sus asíntotas verticales y horizontales.
\end{enumerate}
\end{ejercicios}

\ifmostrarSoluciones
\begin{solucion}[de los ejercicios]
\begin{enumerate}
  \item
  \[
    \lim_{x\to-3^-}\frac{1}{x+3}=-\infty,
    \qquad
    \lim_{x\to-3^+}\frac{1}{x+3}=+\infty.
  \]
  La asíntota vertical es $x=-3$.
  \item El denominador se aproxima a $0$ mediante valores positivos desde ambos lados; por tanto,
  \[
    \lim_{x\to1}\frac{1}{(x-1)^2}=+\infty.
  \]
  La gráfica crece sin cota a ambos lados de la asíntota vertical $x=1$.
  \item No. Para $x\neq2$ se cumple $g(x)=x+2$ y el límite en $2$ es $4$.
  \item
  \begin{enumerate}
    \item Los límites hacia $+\infty$ y $-\infty$ son $0$; la asíntota horizontal es $y=0$.
    \item Los dos límites son $4/2=2$; la asíntota horizontal es $y=2$.
    \item No existe una asíntota horizontal. La función se comporta como $x$; por tanto, el límite hacia $+\infty$ es $+\infty$ y el límite hacia $-\infty$ es $-\infty$.
  \end{enumerate}
  \item El término dominante de $p$ es $-3x^5$. Como la potencia es impar,
  \[
    \lim_{x\to+\infty}p(x)=-\infty,
    \qquad
    \lim_{x\to-\infty}p(x)=+\infty.
  \]
  La función $r$ se comporta como $-x^2$; por tanto,
  \[
    \lim_{x\to+\infty}r(x)
    =\lim_{x\to-\infty}r(x)
    =-\infty.
  \]
  \item Como $\sqrt{x^2+4}=|x|\sqrt{1+4/x^2}$, los límites son
  \[
    \lim_{x\to+\infty}\frac{2x+1}{\sqrt{x^2+4}}=2,
    \qquad
    \lim_{x\to-\infty}\frac{2x+1}{\sqrt{x^2+4}}=-2.
  \]
  \item La asíntota vertical es $x=2$. La información dada garantiza la asíntota horizontal $y=3$ hacia $+\infty$; no se dispone de información sobre el comportamiento cuando $x\to-\infty$.
\end{enumerate}
\end{solucion}
\fi

\begin{nota}[ideas que deben quedar claras]
\begin{itemize}
  \item Una asíntota vertical se decide mediante límites cerca de un número; una asíntota horizontal se decide mediante límites cuando $x\to\pm\infty$.
  \item Un cero del denominador no garantiza una asíntota vertical: primero se revisan factores comunes.
  \item La comparación de grados permite decidir rápidamente los límites al infinito de una función racional.
  \item Cuando domina una potencia, la paridad del exponente y el signo del coeficiente determinan el comportamiento hacia $+\infty$ y $-\infty$.
  \item En expresiones con raíces, debe recordarse que $\sqrt{x^2}=|x|$.
\end{itemize}
\end{nota}

\fuentes{J. Stewart, \emph{Cálculo de una variable: trascendentes tempranas}, apartados sobre límites infinitos, límites al infinito y asíntotas; G. Strang--E. Herman, \emph{Cálculo, volumen 1}, secciones sobre límites al infinito.}

\end{document}
